\documentclass[ekobook.tex]{subfiles}

\begin{document}\setWPP
\setlist[itemize]{
  itemsep=0pt,
  parsep=0pt,      % space between paragraph lines inside one item
  topsep=0.5ex,    % space before/after the whole list (usually keep a little)
  partopsep=0pt
}
\setlength{\columnsep}{0pt}
\lsection{Pojištění}
\qaw{Co je pojištění}{%
Zvláštní odvětví ekonomiky, jedná se o \texthl{finanční službu}, neznamená \texthl{růst bezpečí} ani \texthl{naprostou jistotu}.
}
\qaw{Vyjmenuj činnosti pojišťoven}{%
\begin{itemize}[label={\color{WPP}\faSlackHash}]
\item \textbl{pojišťovací činnost --} uzavírání smluv, vyplácení náhrad, vybírání pojistného, poradenství
\item \textbl{zajišťovací činnost --} zajišťovna pojišťuje pojišťovny
\item \textbl{zábranná činnost --} snaží se předcházet pojistným událostem, osvěta, reklama, školení
\end{itemize}
}
\qaw{Jaké jsou právní formy, kdo uděluje licenci}{%
Nejčastěji se jedná o \texthl{akciovou společnost}, někdy i o \texthl{družstevní organizace}.
Licenci uděluje a spravuje \texthl{ČNB}.
}
\qaw{Vysvětli:}{%
\textbl{pojistník}
-- fyzická nebo právnická osoba, co uzavřela pojistnou smlouvy (např.: rodič na dítě, osoba na auto)\\
\textbl{pojištěný}
-- osoba v jejíž prospěch je smlouva uzavřena\\
\textbl{pojistitel}
-- pojišťovna\\
\textbl{vinkulace}
-- v případě úmrtí pojištěného se z pojistky zaplatí přednostně úvěr (u životního pojištění)\\
\textbl{obmyšlená osoba}
-- osoba, která je uvedena pro vyplacení pojistky (u životního pojištění pozůstalí)\\
\textbl{pojistné plnění}
-- vyplacení pojistného po pojistné události\\
\textbl{výluka z pojištění}
-- pojišťovna zaplatí pouze to, na co je pojistka uzavřena, může být dohodnut i časový limit\\
\textbl{spoluúčast}
-- část škody, kterou pojišťovna nehradí a pojištený si ji tedy hradí sám -- spoluúčastní se na nákladech opravy
}
\vspace*{1em}
\wbreak
\begin{multicols}{2}
\qaw{Jak se pojištění dělí dle formy}{%
\begin{itemize}[label={\color{WPP}\faSlackHash}]
\item \textsc{\color{WPP} Zákonné pojištění}
-- je automaticky uzavřeno dle zákona, \texthl{pro všechny platí stejné podmínky}, je povinné (pojištění zodpovědnosti za provoz motorových vozidel, ZP, SP)\\
\textbl{zdravotní pojištění}\\
-- platí zaměstnanec, zaměstnavatel, OSVČ a stát\\
-- plyne zdravotním pojišťovnám\\
-- hradí se z něj vyšetření a ošetření, částečně léky a pomůcky\\
\textbl{sociální pojištění}\\
-- platí zaměstnanec, zaměstnavatel a OSVČ\\
-- plyne přes ČSSZ do státní poklady\\
-- hradí se z něj podpora v nezaměstnanosti, sociální dávky, důchody, rodičovská dovolená a nemocenská\\
\textbl{pojištění odpovědnosti organizace za škodu způsobenou zdravotním úrazem}\\
-- jedná se o odškodnění pracovních úrazů a nemocí z povolání.
\item \textsc{\color{WPP} Povinně smluvní pojištění}
-- týká se oblastí, kde je \texthl{zvýšené nebezpečí ohrožení ostatních}, patří sem povinné ručení
\item \textsc{\color{WPP} Smluvní pojištění}
-- na základě smlouvy a \texthl{je dobrovolné}, je ovlivněno všeobecnými a zvláštními podmínkami, jednotlivé pojišťovny mohou mít ještě vlastní specifika
\end{itemize}
}
\qaw{Podle formy tvorby rezerv}{%
\begin{itemize}[label={\color{WPP}\faSlackHash}]
\item \textsc{\color{WPP} Rizikové pojištění}
-- neví se kdy nebo jestli pojistná událost vznikne (úrazové, majetkové, ...)
\item \textsc{\color{WPP} Rezervotvorné pojištění}
-- ví se, kdy vznikne (životní a důchodové)
\end{itemize}
}
\qaw{Podle odvětví}{%
\begin{itemize}[label={\color{WPP}\faSlackHash}]
\item \textsc{\color{WPP} Neživotní}
-- na majetek
\item \textsc{\color{WPP} Životní}
-- na osobu
\end{itemize}
}
\qaw{Jaké je pojištění podle druhu}{%
\begin{itemize}[label={\color{WPP}\faSlackHash}]
\item \textsc{\color{WPP} Pojištění majetku a odpovědnosti za škody obyvatelstva}\\
\textbl{pojištění domácnosti}\\
-- jedná se o jedno z nejvýznamnějších pojištění, které využívají asi 2/3 domácností\\
-- kryje převážně živelná rizika, tedy požár; povodeň; sesuv půdy; pády stromů ...
\columnbreak\\
\textbl{pojištění staveb (budovy)}\\
-- týká se rodinných domků, nájemních a obytných domů, zemědělských usedlostí ...\\
\textbl{havarijní pojištění motorových vozidel}\\
-- předmětem jsou motorová vozidla a kryje riziko živelné pohromy, vandalismu, odcizení a havárie při střetu se zvěří
\item \textsc{\color{WPP} Pojištění osob}\\
\textbl{úrazové pojištění}\\
-- finanční zabezpečené pojištěné osoby v případě, že dojde k dočasnému či trvalému tělesnému poškození nebo smrti\\
\textbl{životní pojištění}\\
-- riziko smrti a riziko dožití sjednaného věku\\
\textbl{důchodové pojištění}\\
-- doplňuje důchody poskytované jako dávky sociálního pojištění\\
\textbl{pojištění léčebných výloh v zahraničí}\\
-- úhrada z lékařského hlediska nezbytných nákladů na ošetření
\item \textsc{\color{WPP} Pojištění podnikatelských, průmyslových a zemědělských rizik}\\
\textbl{živelní pojištění}\\
-- nejrozšířenější pojištění právnických osob, týká se movitých i nemovitých věcí ve vlastnictví podnikatelského subjektu\\
\textbl{pojištění krádeží vloupáním nebo loupežným přepadením}\\
-- předmětem jsou movité věci vlastněné podnikatelským subjektem, pojistná částka je odvozena od nové nabo časové hodnoty\\
\textbl{dopravní pojištění}\\
-- vztahuje se na přepravované věci\\
\textbl{úvěrové pojištění}\\
-- kryje před neschopností splácet (dočasný propad peněžitého majetku) a chrání před důsledky nesplacení úvěru\\
\textbl{pojištění zemědělských rizik}\\
-- pojišťuje plodiny a zemědělská zvířata
\end{itemize}
}
\end{multicols}
\vspace*{-1em}
\wbreak
\qaw{Co je malus \wx{ } Co je bonus}{%
Systém bonusů a přirážek používaný u \texthl{havarijního pojištění}. Za období \texthl{bez nehody} dostáváte \texthl{bonusy} (odečty od ceny pojištění). Za nehody se bonusy odečítají až do záporu, kdy vzniká \texthl{riziková přirážka} – tedy \texthl{malus}.
}

\wpage
\end{document}